\section{The engineered full-batch simulation}\label{app:engineered}
The model used by \citet[Theorem~3d]{AKMSS21} is a directed acyclic neural network with $n+1$ inputs, including a constant input, one output, and one trainable weight per edge. Its activation is the piecewise-linear ``two-stage ramp''
\[
 \sigma(z)=\begin{cases}
 -2,&z<-3,\\ z+1,&-3\le z\le-1,\\0,&-1<z<0,\\z,&0\le z\le2,\\2,&z>2.
 \end{cases}
\]
Initialization is a specified polynomial-time computable map $W:\{0,1\}^r\to\R^p$, rather than an independent Gaussian law. Labels lie in $\{0,1\}$ and the loss is $\tfrac12(f(x)-y)^2$. Full-batch updates reuse one iid sample $S$ of size $m_s$, with a fixed constant step size and
\begin{equation}\label{eq:engineered-gradient}
 \theta_{t+1}=\theta_t-\gamma g_t,\qquad
 g_t\in\lambda\mathbb Z^p\cap[-1,1]^p,\qquad
 \left\|g_t-\frac1{m_s}\sum_{(x,y)\in S}
       [\nabla_\theta\tfrac12(f_{\theta_t}(x)-y)^2]_1\right\|_\infty
 \le\frac{3\lambda}{4}.
\end{equation}
Here $[\cdot]_1$ denotes coordinatewise clipping to $[-1,1]$. The guarantee is $\E_{S,R}\sup_{(g_t)\ \mathrm{admissible}}L(f_{\theta_{T\prime}})$: expectation over the reused sample and initialization follows the supremum over all admissible rounded-gradient sequences. The supremum may depend on that entire sample. The constructed trajectories are differentiable and have gradient coordinates bounded by one, so clipping leaves them unchanged.

The imported simulation states that an SQ algorithm with $k$ queries, tolerance $\tau_0$, $r$ random bits and running time $\mathrm{TIME}$ has such a simulation with
\begin{equation}\label{eq:engineered-parameters}
 \lambda=\tau_0/16,\quad T'=2k,\quad
 m_s\lambda^2>C\bigl(k\log(1/\lambda)+\log(1/\delta)\bigr),\quad
 p=\operatorname{poly}(k,1/\tau_0,r,\mathrm{TIME},1/\delta),
\end{equation}
whose expected square loss exceeds the SQ algorithm's by at most $\delta$. The constant $C$ is universal. These are the parameters of the full-batch part of that theorem; the sample is not refreshed at each step.


\subsection{Application to the succinct learner}\label{proof:engineered}
\begin{proof}[Proof of \cref{cor:engineered}]
Use the deterministic improper order-median learner of \cref{thm:fast} at accuracy $\epsilon/8$. It works for every key with $k=O(n^2)$ queries, tolerance $\Omega(\epsilon)$, and polynomial running time including query evaluation. Recode its labels and predictions from $\{-1,+1\}$ to $\{0,1\}$. Its square loss is half its classification loss, hence at most $\epsilon/16$ against every valid response sequence. A dyadic tolerance $\tau_0$ within a factor two of the advertised one preserves this guarantee and the asymptotic bounds.

Give the deterministic source learner one ignored random bit and apply \eqref{eq:engineered-parameters} with $r=1$ and $\delta=\epsilon/16$. Its initialization map may ignore this bit as well. The sample size and model size are polynomial in $n$ for fixed $c,\epsilon$, and the expected final square loss is at most $\epsilon/8$. A thresholding mistake implies $|f(x)-y|\ge1/2$ and therefore square loss at least $1/8$. The expected classification loss is at most $\epsilon$. The PRF sign-rank conclusion is \cref{thm:prf}; it is independent of the simulation. The deterministic source learner satisfies the simulation's supremum-over-responses premise without exchanging a supremum and an expectation.
\end{proof}